\section{Method}
\label{sec:method}

\subsection{Task Setup}

We study question-level examples from ProofWriter. Each example contains a natural-language theory and a query. Under the benchmark's open-world semantics, the query falls into one of three cases:
\begin{itemize}
\item \textbf{True}: the query literal is derivable.
\item \textbf{False}: the opposite literal is derivable.
\item \textbf{Unknown}: neither the query nor its opposite is derivable from the theory.
\end{itemize}
These cases require different responses. A faithful system should prove true queries, refute false queries, and abstain on unsupported ones.

\subsection{Target Format}

We use one structured output format for all three cases:
\begin{verbatim}
MODE: <PROVE|REFUTE|ABSTAIN>
LABEL: <True|False|Unknown>
CHAIN: <triple/rule chain or FAIL[...]>
WITNESS: <statement or missing support>
\end{verbatim}

The meaning of the fields depends on the case:
\begin{itemize}
\item \prove{}: the chain derives the queried literal.
\item \refute{}: the chain derives the opposite literal, which is then stated in the witness.
\item \abstain{}: the chain records a failed-rule tag or a no-rule tag, and the witness states what support is missing.
\end{itemize}

The format is deliberately small. It is easy to generate, easy to parse, and easy to verify. Across variants, we keep the prompt scaffold fixed and change only the supervision target.

\subsection{Building Targets}

We derive question-level targets from the structured ProofWriter data. Facts and rules are already indexed (for example, \texttt{triple8} and \texttt{rule4}), and the benchmark provides proof traces for many true and false queries. We reuse those traces to build \prove{} and \refute{} targets.

Unknown queries need extra processing. We run a lightweight forward-chaining reasoner over the theory and inspect the shallowest blocked derivation for the query. The resulting abstention target is either:
\begin{itemize}
\item a \texttt{FAIL[ruleX]} tag with the missing literals that block that rule, or
\item \texttt{FAIL[no\_rule]} when no rule or fact can support the query.
\end{itemize}
This turns an unsupported query into a concrete, checkable target instead of only the label \texttt{Unknown}.

\subsection{Verifier and Metrics}

We evaluate each output at two levels:
\begin{itemize}
\item \textbf{Label correctness}: whether the predicted label matches the benchmark answer.
\item \textbf{Evidence faithfulness}: whether the generated chain or abstention witness is symbolically consistent with the theory.
\end{itemize}
For \prove{} and \refute{} outputs, we replay the predicted fact/rule chain and check that it derives the claimed literal. For \abstain{} outputs, we check that the stated missing support is absent from the theory closure. We report \emph{joint accuracy}, which requires both a correct label and faithful evidence.

The verifier is intentionally lightweight. It does not search for a proof on the model's behalf or repair an incorrect output. It only checks the evidence the model chose to provide.

\subsection{Scope of the Method}

\model{} changes the target format rather than the inference stack. There is no planner, no search controller, and no external solver in the main test-time loop. This keeps the empirical question narrow: can better supervision alone improve verifier-backed proof, refutation, and abstention behavior?
